% chapters/02/checklist.tex
\bagian{Daftar Periksa Pemahaman Konsep}

Bab ini tidak menuntut keterampilan menulis program, melainkan pemahaman
konsep. Gunakan daftar berikut untuk menilai kesiapan Anda memasuki Bab~III.
Butir yang belum tercentang menandakan bagian yang perlu Anda baca ulang.

\begin{checklistbox}
\begin{checklist}
  \item Saya dapat menyebutkan empat unsur pembentuk permainan dan
        menerapkannya pada satu permainan yang saya kenal.
  \item Saya dapat menjelaskan mengapa program permainan disebut sistem waktu
        nyata.
  \item Saya dapat menyebutkan ketiga tahap siklus permainan beserta isi
        pekerjaan masing-masing.
  \item Saya dapat menjelaskan mengapa kecepatan harus dinyatakan dalam satuan
        per detik, bukan per bingkai.
  \item Saya dapat menghitung perpindahan objek pada satu bingkai dari
        kecepatan dan selang waktunya.
  \item Saya dapat menjelaskan perbedaan langkah tetap dan langkah berubah
        beserta kelebihan masing-masing.
  \item Saya dapat menjelaskan arah sumbu pada koordinat layar dan akibatnya
        pada penulisan gravitasi.
  \item Saya dapat menjelaskan apa yang dilakukan persegi pembatas dan apa
        kekurangan pendekatan tersebut.
  \item Saya dapat menjelaskan mengapa galat kaidah Euler bergantung pada
        panjang langkah.
  \item Saya dapat membedakan lapisan keputusan dan lapisan navigasi pada
        pengendalian NPC.
  \item Saya dapat menjelaskan perbedaan cara kerja pustaka dan mesin
        permainan.
  \item Saya dapat menyebutkan padanan konsep antara Python/Pygame dan Unity
        untuk sekurang-kurangnya lima konsep pada
        Tabel~\ref{tab:padanan-konsep}.
\end{checklist}
\end{checklistbox}
